\subsection{Macroscopic chords in a hexagon}\label{sec:pv:hexagon-proof}

For specificity take regular lattice hexagons of integer side \(R\to\infty\). Say an arc is macroscopic when its diameter is at least \(R/10\). Their absolute critical mass over all boundary sources is
\[
 R^{3/4+o(1)}.
\]
For the lower bound use sources in the central third of one side. Exits not to the initial bottom side all have macroscopic reach. Flux and comparison of bottom returns with arches in a containing flat strip of \(O(R)\) layers give non-bottom mass at least \(b_{O(R)}\).

For the upper bound per source, impose a cut parallel to its starting side at height $h=\lfloor R/\log^2R\rfloor$. Flux subtraction bounds the mass of exits removed by this cut by a constant times $b_h$, since its new exit prefixes are strip bridges. A retained macroscopic path must reach horizontal distance comparable to $R$, so Lemma~\ref{lem:pv:lateral-tail} bounds its weight by $Ce^{-c\log^2R}$. Summing over the $O(R)$ boundary sources proves the upper normalizing exponent.

Under the normalized macroscopic arc weights, bulk vertices (say in the central third) now have visit probability
\[
 R^{-2/3+o(1)}.
\]
Indeed all visiting paths there qualify, and the surrounding-polygon sums in Proposition~\ref{prop:pv:visits} lie between disk sums at comparable radii. Proposition~\ref{prop:pv:planar-localization} gives their exponent. The same upper bound holds throughout the hexagon. Summing the probabilities yields the mean visited mass, or length,
\[
 \mathbb E |\gamma|=R^{4/3+o(1)} .
\]

\subsection{A fixed-source strip version}\label{c01-c1-a-fixed-source-strip-version}

There is also an annealed dimension statement with one actual boundary point fixed. In the full planar strip of \(h\) layers take all the critical exits from a fixed bottom port (thus bridges to the top \textbf{or} arches back to bottom), conditioned to get at least one quarter of the way up. Then
\[
 \mathbb E |\gamma|=h^{4/3+o(1)} .
\]
Changing this to any fixed positive fraction strictly below one half gives the same result. The exit mass of this ensemble before normalization is \(h^{-1/4+o(1)}\): lower by the bridges, upper for the arches by the height-tail bound. This conditioning could also confine horizontal range to \(h\log^2 h\) about the start without changing the exponent. For bridges alone the argument gives the mean upper bound \(h^{4/3+o(1)}\) (and hence the usual in-probability upper-exponent bound), not by itself the matching bridge lower bound.

The pointwise estimate behind this statement will also be needed for the
bridge lower bound. It sums over all boundary positions, rather than
fixing a source.

\begin{lemma}[Strip visit weights]\label{lem:pv:strip-visits}
Let $A_{S_h}(v)$ be the positive weight of ordered strip excursions
through $v$, with both endpoints free on the two horizontal boundaries.
Uniformly over strip vertices,
\[
 A_{S_h}(v)\le h^{1/12+o(1)}.
\]
For every fixed $\delta>0$, uniformly over vertices at least
$\delta h$ layers from both horizontal boundaries,
\[
 A_{S_h}(v)=h^{1/12+o(1)}.
\]
The same upper bound holds for the ordered excursion weight $A_D(v)$
in every convex lateral truncation $D$, including excursions with
lateral endpoints. Its lower bound holds uniformly when, in addition
to the horizontal clearance above, $v$ is at distance at least
$\delta h$ from every lateral side.
\end{lemma}

\begin{proof}
Convexly truncate the strip at lateral distance $W$ on both sides using
tiling lines of the two slanted side orientations. For a face within
bounded distance of the strip, fugacity-$2$ systems using strip vertices
and surrounding that face have sum at most $h^{1/12+o(1)}$, uniformly in
the truncation. Indeed every enclosing polygon interior lies in the
slab and is disjoint from its translates by nonzero multiples of
$p\omega$, for a sufficiently large integral $p\asymp h$.
Corollary~\ref{cor:pv:oversized-polygons} applies to this transverse
period: its geometric argument rotates the period direction to
horizontal, while its degree bound uses a sufficiently large balanced
cylinder containing any finite test subset and all its rotations. Loops of
radius greater than $p^{1+\alpha}$ therefore cost at most a bounded
exponential factor. The remaining loops are controlled by
Proposition~\ref{prop:pv:planar-localization}. First let $h\to\infty$
for fixed $\alpha>0$, then let
$\alpha\downarrow0$. For faces at least a fixed multiple of $h$ from
the walls and the lateral ends, an inscribed disk gives the matching lower
exponent.

Proposition~\ref{prop:pv:visits} now gives both estimates in every
finite convex truncation. We also need a tail bound for its side
sources. Use the two slanted tiling cuts facing toward the far part of
the strip, with both outward normals pointing laterally that way.
Their distinct angles from the initial inward direction lie strictly
between $-\pi$ and $\pi$, so the separated-ray argument of
Lemma~\ref{lem:pv:lateral-tail} bounds the weight reaching lateral
distance $t$ by $Ce^{-ct/(h+1)}$. Displacement toward the near end is
limited by $O(h+1)$, and length is bounded by
$C(h+1)(1+\text{lateral span})$. Thus each side source has bounded
length mass at fixed $h$, uniformly in the truncation width.

Let $W\to\infty$ at fixed $h$. Excursions between horizontal boundary
ports exhaust the strip excursions. For a fixed marked vertex, paths
from the $O(h)$ distant side ports have vanishing weight by the same
tail bound. This proves the stated full-strip estimates.
\end{proof}

We finish Proposition~\ref{prop:pv:mixed} by averaging over horizontal
sources. In a finite truncation $D$, counting every visited vertex gives
the exact identity
\[
 \sum_{v\in D} A_D(v)
 =\sum_{\gamma\text{ an ordered excursion in }D}
   |\gamma|x_*^{|\gamma|}.
\]
Divide by horizontal width and let that width tend to infinity at fixed
$h$. Lemma~\ref{lem:pv:strip-visits} bounds the result above by
$h^{13/12+o(1)}$. A central band of $\asymp h$ layers supplies the same
lower exponent. For the ensemble conditioned to reach $\vartheta h$
layers above its starting wall, where $0<\vartheta<1/2$ is fixed,
choose this band symmetric about the midline and more than
$\vartheta h$ layers from either wall. Every
path contributing a visit to that band then meets the conditioning
event, from either horizontal starting side.

The slanted side sources disappear in this average: there are only
$O(h)$ of them, and the tail estimate in the preceding proof bounds
each one's length mass at fixed $h$. Their contribution divided by
width therefore vanishes.

Each horizontal wall has one source port per unit length. Translation
invariance and top--bottom symmetry therefore identify the limiting
average with twice the fixed-source length weight. One can first
restrict to bounded path range, pass to
the large-width limit, and then remove that restriction using the
exponential tails. Thus the conditioned ensemble has unnormalized
length weight $h^{13/12+o(1)}$. Its mass is $h^{-1/4+o(1)}$ by
Proposition~\ref{prop:pv:boundary}, giving the asserted mean. The same
tails make the discarded first moment at horizontal range
$h\log^2 h$ smaller than every fixed power of $h$, proving the optional
confinement statement.
